% =====================================================================
%  elegant-style.tex —— PhotoLib 使用指南三份文档共用的样式与宏定义
%  由 管理员/部长/校区负责人 三份指南通过 % =====================================================================
%  elegant-style.tex —— PhotoLib 使用指南三份文档共用的样式与宏定义
%  由 管理员/部长/校区负责人 三份指南通过 % =====================================================================
%  elegant-style.tex —— PhotoLib 使用指南三份文档共用的样式与宏定义
%  由 管理员/部长/校区负责人 三份指南通过 % =====================================================================
%  elegant-style.tex —— PhotoLib 使用指南三份文档共用的样式与宏定义
%  由 管理员/部长/校区负责人 三份指南通过 \input{elegant-style} 引入。
%  依赖：ElegantBook 文档类（须以 XeLaTeX 编译）。
% =====================================================================

% ---- 颜色（与 ElegantBook green 主题协调）----
\definecolor{plMain}{RGB}{0,120,80}      % 主色：墨绿
\definecolor{plInk}{RGB}{38,50,56}       % 正文强调
\definecolor{admTipBg}{RGB}{236,247,240}
\definecolor{admTipFr}{RGB}{34,139,88}
\definecolor{admNoteBg}{RGB}{232,240,254}
\definecolor{admNoteFr}{RGB}{25,103,210}
\definecolor{admWarnBg}{RGB}{253,237,237}
\definecolor{admWarnFr}{RGB}{198,40,40}
\definecolor{admStepBg}{RGB}{245,243,255}
\definecolor{admStepFr}{RGB}{103,58,183}

% ---- tcolorbox（ElegantBook 已载入，这里补声明所需库）----
\usepackage{tcolorbox}
\tcbuselibrary{skins,breakable}

% 统一的四种“提示框”环境，均支持可选自定义标题：
%   \begin{tipbox}[自定义标题] ... \end{tipbox}
\newtcolorbox{tipbox}[1][小贴士]{%
  enhanced, breakable,
  colback=admTipBg, colframe=admTipFr, coltitle=white,
  fonttitle=\bfseries\sffamily, title={#1},
  boxrule=0.6pt, arc=3pt, left=8pt, right=8pt, top=6pt, bottom=6pt}

\newtcolorbox{notebox}[1][说明]{%
  enhanced, breakable,
  colback=admNoteBg, colframe=admNoteFr, coltitle=white,
  fonttitle=\bfseries\sffamily, title={#1},
  boxrule=0.6pt, arc=3pt, left=8pt, right=8pt, top=6pt, bottom=6pt}

\newtcolorbox{warnbox}[1][注意]{%
  enhanced, breakable,
  colback=admWarnBg, colframe=admWarnFr, coltitle=white,
  fonttitle=\bfseries\sffamily, title={#1},
  boxrule=0.6pt, arc=3pt, left=8pt, right=8pt, top=6pt, bottom=6pt}

\newtcolorbox{stepbox}[1][操作步骤]{%
  enhanced, breakable,
  colback=admStepBg, colframe=admStepFr, coltitle=white,
  fonttitle=\bfseries\sffamily, title={#1},
  boxrule=0.6pt, arc=3pt, left=8pt, right=8pt, top=6pt, bottom=6pt}

% ---- 表格 ----
\usepackage{booktabs}
\usepackage{array}
\usepackage{longtable}
\usepackage{makecell}
\renewcommand{\arraystretch}{1.28}

% 左对齐、可换行的定宽列
\newcolumntype{L}[1]{>{\raggedright\arraybackslash}p{#1}}

% ---- 行内“键位/菜单/字段”标记 ----
\usepackage{relsize}
% 菜单路径：系统管理 \menu{数据备份}
\newcommand{\menu}[1]{{\sffamily\bfseries\color{plMain}#1}}
% 界面按钮/字段名
\newcommand{\ui}[1]{\textsf{「#1」}}
% 权限码等等宽标识；在每个下划线后允许换行，避免长标识顶出表格边界
\newcommand{\code}[1]{{\ttfamily\smaller[0.5]%
  \renewcommand{\_}{\textunderscore\penalty0\hspace{0pt}}#1}}
% 角色徽标
\newcommand{\role}[1]{{\sffamily\bfseries\color{admNoteFr}#1}}

% ---- 超链接 ----
\hypersetup{
  colorlinks=true,
  linkcolor=plMain,
  urlcolor=admNoteFr,
  citecolor=plMain}
 引入。
%  依赖：ElegantBook 文档类（须以 XeLaTeX 编译）。
% =====================================================================

% ---- 颜色（与 ElegantBook green 主题协调）----
\definecolor{plMain}{RGB}{0,120,80}      % 主色：墨绿
\definecolor{plInk}{RGB}{38,50,56}       % 正文强调
\definecolor{admTipBg}{RGB}{236,247,240}
\definecolor{admTipFr}{RGB}{34,139,88}
\definecolor{admNoteBg}{RGB}{232,240,254}
\definecolor{admNoteFr}{RGB}{25,103,210}
\definecolor{admWarnBg}{RGB}{253,237,237}
\definecolor{admWarnFr}{RGB}{198,40,40}
\definecolor{admStepBg}{RGB}{245,243,255}
\definecolor{admStepFr}{RGB}{103,58,183}

% ---- tcolorbox（ElegantBook 已载入，这里补声明所需库）----
\usepackage{tcolorbox}
\tcbuselibrary{skins,breakable}

% 统一的四种“提示框”环境，均支持可选自定义标题：
%   \begin{tipbox}[自定义标题] ... \end{tipbox}
\newtcolorbox{tipbox}[1][小贴士]{%
  enhanced, breakable,
  colback=admTipBg, colframe=admTipFr, coltitle=white,
  fonttitle=\bfseries\sffamily, title={#1},
  boxrule=0.6pt, arc=3pt, left=8pt, right=8pt, top=6pt, bottom=6pt}

\newtcolorbox{notebox}[1][说明]{%
  enhanced, breakable,
  colback=admNoteBg, colframe=admNoteFr, coltitle=white,
  fonttitle=\bfseries\sffamily, title={#1},
  boxrule=0.6pt, arc=3pt, left=8pt, right=8pt, top=6pt, bottom=6pt}

\newtcolorbox{warnbox}[1][注意]{%
  enhanced, breakable,
  colback=admWarnBg, colframe=admWarnFr, coltitle=white,
  fonttitle=\bfseries\sffamily, title={#1},
  boxrule=0.6pt, arc=3pt, left=8pt, right=8pt, top=6pt, bottom=6pt}

\newtcolorbox{stepbox}[1][操作步骤]{%
  enhanced, breakable,
  colback=admStepBg, colframe=admStepFr, coltitle=white,
  fonttitle=\bfseries\sffamily, title={#1},
  boxrule=0.6pt, arc=3pt, left=8pt, right=8pt, top=6pt, bottom=6pt}

% ---- 表格 ----
\usepackage{booktabs}
\usepackage{array}
\usepackage{longtable}
\usepackage{makecell}
\renewcommand{\arraystretch}{1.28}

% 左对齐、可换行的定宽列
\newcolumntype{L}[1]{>{\raggedright\arraybackslash}p{#1}}

% ---- 行内“键位/菜单/字段”标记 ----
\usepackage{relsize}
% 菜单路径：系统管理 \menu{数据备份}
\newcommand{\menu}[1]{{\sffamily\bfseries\color{plMain}#1}}
% 界面按钮/字段名
\newcommand{\ui}[1]{\textsf{「#1」}}
% 权限码等等宽标识；在每个下划线后允许换行，避免长标识顶出表格边界
\newcommand{\code}[1]{{\ttfamily\smaller[0.5]%
  \renewcommand{\_}{\textunderscore\penalty0\hspace{0pt}}#1}}
% 角色徽标
\newcommand{\role}[1]{{\sffamily\bfseries\color{admNoteFr}#1}}

% ---- 超链接 ----
\hypersetup{
  colorlinks=true,
  linkcolor=plMain,
  urlcolor=admNoteFr,
  citecolor=plMain}
 引入。
%  依赖：ElegantBook 文档类（须以 XeLaTeX 编译）。
% =====================================================================

% ---- 颜色（与 ElegantBook green 主题协调）----
\definecolor{plMain}{RGB}{0,120,80}      % 主色：墨绿
\definecolor{plInk}{RGB}{38,50,56}       % 正文强调
\definecolor{admTipBg}{RGB}{236,247,240}
\definecolor{admTipFr}{RGB}{34,139,88}
\definecolor{admNoteBg}{RGB}{232,240,254}
\definecolor{admNoteFr}{RGB}{25,103,210}
\definecolor{admWarnBg}{RGB}{253,237,237}
\definecolor{admWarnFr}{RGB}{198,40,40}
\definecolor{admStepBg}{RGB}{245,243,255}
\definecolor{admStepFr}{RGB}{103,58,183}

% ---- tcolorbox（ElegantBook 已载入，这里补声明所需库）----
\usepackage{tcolorbox}
\tcbuselibrary{skins,breakable}

% 统一的四种“提示框”环境，均支持可选自定义标题：
%   \begin{tipbox}[自定义标题] ... \end{tipbox}
\newtcolorbox{tipbox}[1][小贴士]{%
  enhanced, breakable,
  colback=admTipBg, colframe=admTipFr, coltitle=white,
  fonttitle=\bfseries\sffamily, title={#1},
  boxrule=0.6pt, arc=3pt, left=8pt, right=8pt, top=6pt, bottom=6pt}

\newtcolorbox{notebox}[1][说明]{%
  enhanced, breakable,
  colback=admNoteBg, colframe=admNoteFr, coltitle=white,
  fonttitle=\bfseries\sffamily, title={#1},
  boxrule=0.6pt, arc=3pt, left=8pt, right=8pt, top=6pt, bottom=6pt}

\newtcolorbox{warnbox}[1][注意]{%
  enhanced, breakable,
  colback=admWarnBg, colframe=admWarnFr, coltitle=white,
  fonttitle=\bfseries\sffamily, title={#1},
  boxrule=0.6pt, arc=3pt, left=8pt, right=8pt, top=6pt, bottom=6pt}

\newtcolorbox{stepbox}[1][操作步骤]{%
  enhanced, breakable,
  colback=admStepBg, colframe=admStepFr, coltitle=white,
  fonttitle=\bfseries\sffamily, title={#1},
  boxrule=0.6pt, arc=3pt, left=8pt, right=8pt, top=6pt, bottom=6pt}

% ---- 表格 ----
\usepackage{booktabs}
\usepackage{array}
\usepackage{longtable}
\usepackage{makecell}
\renewcommand{\arraystretch}{1.28}

% 左对齐、可换行的定宽列
\newcolumntype{L}[1]{>{\raggedright\arraybackslash}p{#1}}

% ---- 行内“键位/菜单/字段”标记 ----
\usepackage{relsize}
% 菜单路径：系统管理 \menu{数据备份}
\newcommand{\menu}[1]{{\sffamily\bfseries\color{plMain}#1}}
% 界面按钮/字段名
\newcommand{\ui}[1]{\textsf{「#1」}}
% 权限码等等宽标识；在每个下划线后允许换行，避免长标识顶出表格边界
\newcommand{\code}[1]{{\ttfamily\smaller[0.5]%
  \renewcommand{\_}{\textunderscore\penalty0\hspace{0pt}}#1}}
% 角色徽标
\newcommand{\role}[1]{{\sffamily\bfseries\color{admNoteFr}#1}}

% ---- 超链接 ----
\hypersetup{
  colorlinks=true,
  linkcolor=plMain,
  urlcolor=admNoteFr,
  citecolor=plMain}
 引入。
%  依赖：ElegantBook 文档类（须以 XeLaTeX 编译）。
% =====================================================================

% ---- 颜色（与 ElegantBook green 主题协调）----
\definecolor{plMain}{RGB}{0,120,80}      % 主色：墨绿
\definecolor{plInk}{RGB}{38,50,56}       % 正文强调
\definecolor{admTipBg}{RGB}{236,247,240}
\definecolor{admTipFr}{RGB}{34,139,88}
\definecolor{admNoteBg}{RGB}{232,240,254}
\definecolor{admNoteFr}{RGB}{25,103,210}
\definecolor{admWarnBg}{RGB}{253,237,237}
\definecolor{admWarnFr}{RGB}{198,40,40}
\definecolor{admStepBg}{RGB}{245,243,255}
\definecolor{admStepFr}{RGB}{103,58,183}

% ---- tcolorbox（ElegantBook 已载入，这里补声明所需库）----
\usepackage{tcolorbox}
\tcbuselibrary{skins,breakable}

% 统一的四种“提示框”环境，均支持可选自定义标题：
%   \begin{tipbox}[自定义标题] ... \end{tipbox}
\newtcolorbox{tipbox}[1][小贴士]{%
  enhanced, breakable,
  colback=admTipBg, colframe=admTipFr, coltitle=white,
  fonttitle=\bfseries\sffamily, title={#1},
  boxrule=0.6pt, arc=3pt, left=8pt, right=8pt, top=6pt, bottom=6pt}

\newtcolorbox{notebox}[1][说明]{%
  enhanced, breakable,
  colback=admNoteBg, colframe=admNoteFr, coltitle=white,
  fonttitle=\bfseries\sffamily, title={#1},
  boxrule=0.6pt, arc=3pt, left=8pt, right=8pt, top=6pt, bottom=6pt}

\newtcolorbox{warnbox}[1][注意]{%
  enhanced, breakable,
  colback=admWarnBg, colframe=admWarnFr, coltitle=white,
  fonttitle=\bfseries\sffamily, title={#1},
  boxrule=0.6pt, arc=3pt, left=8pt, right=8pt, top=6pt, bottom=6pt}

\newtcolorbox{stepbox}[1][操作步骤]{%
  enhanced, breakable,
  colback=admStepBg, colframe=admStepFr, coltitle=white,
  fonttitle=\bfseries\sffamily, title={#1},
  boxrule=0.6pt, arc=3pt, left=8pt, right=8pt, top=6pt, bottom=6pt}

% ---- 表格 ----
\usepackage{booktabs}
\usepackage{array}
\usepackage{longtable}
\usepackage{makecell}
\renewcommand{\arraystretch}{1.28}

% 左对齐、可换行的定宽列
\newcolumntype{L}[1]{>{\raggedright\arraybackslash}p{#1}}

% ---- 行内“键位/菜单/字段”标记 ----
\usepackage{relsize}
% 菜单路径：系统管理 \menu{数据备份}
\newcommand{\menu}[1]{{\sffamily\bfseries\color{plMain}#1}}
% 界面按钮/字段名
\newcommand{\ui}[1]{\textsf{「#1」}}
% 权限码等等宽标识；在每个下划线后允许换行，避免长标识顶出表格边界
\newcommand{\code}[1]{{\ttfamily\smaller[0.5]%
  \renewcommand{\_}{\textunderscore\penalty0\hspace{0pt}}#1}}
% 角色徽标
\newcommand{\role}[1]{{\sffamily\bfseries\color{admNoteFr}#1}}

% ---- 超链接 ----
\hypersetup{
  colorlinks=true,
  linkcolor=plMain,
  urlcolor=admNoteFr,
  citecolor=plMain}
